%======================================================================
%  COMPUTATIONS.TEX -- Sections 7--8 of the standalone paper
%  The genus-2 free energy of W_3 and the gravitational coproduct
%======================================================================


%======================================================================
\section{The genus-2 free energy of \texorpdfstring{$W_3$}{W3}}
\label{sec:genus2-w3}
%======================================================================

The preceding sections developed the theory in the single-channel
regime: one strong generator, one primary line, one modular
characteristic~$\kappa$.  We now compute the first explicit
multi-channel invariant.  The $W_3$ algebra has two strong
generators---the stress tensor~$T$ of conformal weight~$2$ and the
spin-3 current~$W$ of conformal weight~$3$---and the genus-2 free
energy $F_2(W_3)$ receives contributions from \emph{both} channels
and their interactions.

\subsection{The Frobenius algebra of the bar-complex CohFT}
\label{subsec:w3-frobenius}

The bar-complex CohFT of $W_3$ has a two-dimensional underlying
Frobenius algebra with basis $\{e_T, e_W\}$.  The metric is diagonal,
determined by the per-channel modular characteristics:
\begin{equation}\label{eq:w3-metric}
\eta_{TT} = \kappa_T = \frac{c}{2},
\qquad
\eta_{WW} = \kappa_W = \frac{c}{3},
\qquad
\eta_{TW} = 0.
\end{equation}
The total modular characteristic is $\kappa(W_3) = \kappa_T + \kappa_W
= 5c/6 = c \cdot (H_3 - 1)$, where $H_3 = 1 + 1/2 + 1/3 = 11/6$
is the third harmonic number.  The inverse metric gives the
propagators:
\begin{equation}\label{eq:w3-propagators}
\eta^{TT} = \frac{2}{c},
\qquad
\eta^{WW} = \frac{3}{c}.
\end{equation}
Both propagators have weight~$1$:
the bar-complex propagator is $d\!\log E(z,w)$, where
$E(z,w)$ is the prime form (a section of
$K^{-1/2} \boxtimes K^{-1/2}$), so $d\!\log E$ has weight~$1$
regardless of the conformal weight of the field being sewed.

The three-point structure constants are determined by the OPE.
The $\mathbb{Z}_2$ symmetry $W \mapsto -W$ forces all three-point
functions with an odd number of $W$-insertions to vanish:
\begin{equation}\label{eq:w3-C3}
C_{TTT} = C_{TWW} = C_{WWT} = c,
\qquad
C_{TTW} = C_{TWT} = C_{WTT} = C_{WWW} = 0.
\end{equation}
The derivations are:
$T_{(1)}T = 2T$ gives $C^T_{TT} = 2$, whence
$C_{TTT} = \eta_{TT} \cdot 2 = c$;
$T_{(1)}W = 3W$ gives $C^W_{TW} = 3$, whence
$C_{TWW} = \eta_{WW} \cdot 3 = c$;
$W_{(3)}W = 2T$ gives $C^T_{WW} = 2$, whence
$C_{WWT} = \eta_{TT} \cdot 2 = c$.

The genus-0 four-point vertex exhibits a remarkable universality.
For any channel pair $(i,j) \in \{T,W\}^2$:
\begin{equation}\label{eq:w3-V04}
V_{0,4}(i,i \mid j,j)
\;=\;
\sum_m \eta^{mm}\, C_{iim}\, C_{jjm}
\;=\;
\eta^{TT} \cdot c \cdot c
\;=\;
2c.
\end{equation}
Only the $T$-channel contributes to the sum: $C_{iiT} = c$ for
$i \in \{T,W\}$, while $C_{iiW} = 0$ for both $i = T$ (since
$C_{TTW} = 0$) and $i = W$ (since $C_{WWW} = 0$).
The $W$-channel is invisible at genus-0 valence~4.


\subsection{The seven stable graphs at genus 2}
\label{subsec:genus2-graphs}

The moduli space $\overline{\mathcal{M}}_{2,0}$ has a
stratification by seven stable graphs $\Gamma_0, \ldots, \Gamma_6$:

\begin{center}
\small
\renewcommand{\arraystretch}{1.3}
\begin{tabular}{clllccl}
& \textbf{Name}
& \textbf{Vertices}
& \textbf{Edges}
& $|\mathrm{Aut}|$
& $h^1$
& \textbf{Type} \\
\hline
$\Gamma_0$ & smooth
  & $(g{=}2,\, \mathrm{val}{=}0)$
  & $0$
  & $1$ & $0$ & compact stratum \\[2pt]
$\Gamma_1$ & figure-eight
  & $(g{=}1,\, \mathrm{val}{=}2)$
  & $1$ self-loop
  & $2$ & $1$ & one-loop \\[2pt]
$\Gamma_2$ & banana
  & $(g{=}0,\, \mathrm{val}{=}4)$
  & $2$ self-loops
  & $8$ & $2$ & two-loop sunset \\[2pt]
$\Gamma_3$ & dumbbell
  & $(g{=}1,\, \mathrm{val}{=}1) + (g{=}1,\, \mathrm{val}{=}1)$
  & $1$ bridge
  & $2$ & $0$ & separating \\[2pt]
$\Gamma_4$ & theta
  & $(g{=}0,\, \mathrm{val}{=}3) + (g{=}0,\, \mathrm{val}{=}3)$
  & $3$ bridges
  & $12$ & $2$ & non-separating \\[2pt]
$\Gamma_5$ & lollipop
  & $(g{=}0,\, \mathrm{val}{=}3) + (g{=}1,\, \mathrm{val}{=}1)$
  & $1$ self-loop $+$ $1$ bridge
  & $2$ & $1$ & mixed
\end{tabular}
\end{center}
Every vertex satisfies the stability condition $2g_v + \mathrm{val}_v
\geq 3$, and the genus formula $h^1(\Gamma) + \sum_v g_v = 2$ holds
in each case.  Note that the smooth graph $\Gamma_0$ has no edges and
contributes the integrated class on $\mathcal{M}_{2,0}$ itself; the
remaining five graphs are the boundary strata.


\subsection{Multi-channel Feynman rules}
\label{subsec:w3-feynman-rules}

For each graph $\Gamma$ with $|E(\Gamma)|$ edges, a \emph{channel
assignment} is a function $\sigma \colon E(\Gamma) \to \{T, W\}$.
The Feynman rules of the bar-complex CohFT assign to each pair
$(\Gamma, \sigma)$ an amplitude
\begin{equation}\label{eq:feynman-rule}
\mathcal{A}(\Gamma, \sigma)
\;=\;
\frac{1}{|\mathrm{Aut}(\Gamma)|}
\prod_{e \in E(\Gamma)} \eta^{\sigma(e),\, \sigma(e)}
\prod_{v \in V(\Gamma)} V_v\bigl(g_v;\, \sigma|_{\mathrm{half\text{-}edges}(v)}\bigr),
\end{equation}
where the vertex factors are:
\begin{itemize}
\item \emph{Genus-0, trivalent}: $V_v = C_{\sigma(e_1)\sigma(e_2)\sigma(e_3)}$.
\item \emph{Genus-0, four-valent} (the banana vertex):
  $V_v = \sum_m \eta^{mm} C_{\sigma(e_1)\sigma(e_2)m}\,
         C_{m\,\sigma(e_3)\sigma(e_4)}$.
\item \emph{Genus-1, valence $n$}:
  $V_v = \kappa_{\sigma(e)}/24$ (per-channel genus-1 universality,
  proved unconditionally).
\end{itemize}
Each half-edge at a self-loop contributes a pair of half-edges to
the same vertex, both carrying channel~$\sigma(e)$.


\subsection{Per-graph amplitudes}
\label{subsec:per-graph-amplitudes}

We now compute the amplitude of each boundary graph, decomposed into
\emph{diagonal} (all edges carry the same channel) and \emph{mixed}
(at least two edges carry different channels) contributions.

\subsubsection*{$\Gamma_1$: figure-eight}

One edge, two channel assignments.  The vertex is genus-1 with
valence~2, so the vertex factor for channel~$i$ is $\kappa_i/24$:
\begin{align}
\mathcal{A}(\Gamma_1, T) &= \tfrac{1}{2}\, \eta^{TT} \cdot
  \tfrac{\kappa_T}{24}
  = \tfrac{1}{2} \cdot \tfrac{2}{c} \cdot \tfrac{c}{48}
  = \tfrac{1}{48}, \label{eq:fig-eight-T}\\
\mathcal{A}(\Gamma_1, W) &= \tfrac{1}{2}\, \eta^{WW} \cdot
  \tfrac{\kappa_W}{24}
  = \tfrac{1}{2} \cdot \tfrac{3}{c} \cdot \tfrac{c}{72}
  = \tfrac{1}{48}. \label{eq:fig-eight-W}
\end{align}
Both channels contribute equally: the propagator~$1/\kappa_i$ cancels
the genus-1 vertex factor~$\kappa_i/24$, leaving the universal value
$1/48$ independent of $c$.  The total is $1/24 = \lambda_1^{\mathrm{FP}}$.
No mixed channel assignments exist (single edge).

\subsubsection*{$\Gamma_2$: banana}

Two edges, four channel assignments, $|\mathrm{Aut}| = 8$.  The
vertex is genus-0 with valence~4, and the four-point vertex factor
is $V_{0,4}(i,i \mid j,j) = 2c$ universally~\eqref{eq:w3-V04}.
\begin{align}
\mathcal{A}(\Gamma_2, TT) &=
  \tfrac{1}{8}\, (\eta^{TT})^2 \cdot 2c
  = \tfrac{1}{8} \cdot \tfrac{4}{c^2} \cdot 2c
  = \frac{1}{c}, \label{eq:banana-TT}\\
\mathcal{A}(\Gamma_2, WW) &=
  \tfrac{1}{8}\, (\eta^{WW})^2 \cdot 2c
  = \tfrac{1}{8} \cdot \tfrac{9}{c^2} \cdot 2c
  = \frac{9}{4c}, \label{eq:banana-WW}\\
\mathcal{A}(\Gamma_2, TW{+}WT) &=
  \tfrac{1}{8} \cdot 2 \cdot \eta^{TT}\, \eta^{WW} \cdot 2c
  = \tfrac{1}{8} \cdot 2 \cdot \tfrac{6}{c^2} \cdot 2c
  = \frac{3}{c}. \label{eq:banana-mixed}
\end{align}
The cross-channel contribution is $\delta_{\mathrm{banana}} = 3/c$.

\subsubsection*{$\Gamma_3$: dumbbell}

One bridge edge, two channel assignments, $|\mathrm{Aut}| = 2$.
Each vertex has genus~1 and valence~1, with vertex factor
$\kappa_i/24$:
\begin{align}
\mathcal{A}(\Gamma_3, T) &=
  \tfrac{1}{2}\, \eta^{TT} \cdot
  \bigl(\tfrac{\kappa_T}{24}\bigr)^2
  = \tfrac{1}{2} \cdot \tfrac{2}{c} \cdot \tfrac{c^2}{2304}
  = \frac{c}{2304}, \label{eq:dumbbell-T}\\
\mathcal{A}(\Gamma_3, W) &=
  \tfrac{1}{2}\, \eta^{WW} \cdot
  \bigl(\tfrac{\kappa_W}{24}\bigr)^2
  = \tfrac{1}{2} \cdot \tfrac{3}{c} \cdot \tfrac{c^2}{5184}
  = \frac{c}{3456}. \label{eq:dumbbell-W}
\end{align}
The total is $c/2304 + c/3456 = (3c + 2c)/6912 = 5c/6912 =
\kappa(W_3)/1152$.  No mixed channel assignments exist (single edge).

\subsubsection*{$\Gamma_4$: theta}

Three bridge edges, eight channel assignments, $|\mathrm{Aut}| = 12$.
Each vertex has genus~0 and valence~3, with vertex factor
$C_{\sigma(e_1)\sigma(e_2)\sigma(e_3)}$.
\begin{itemize}
\item \emph{All-$T$}: both vertices get $C_{TTT} = c$.
  Amplitude $= \frac{1}{12}\, (\eta^{TT})^3 \cdot c^2
  = \frac{1}{12} \cdot \frac{8}{c^3} \cdot c^2
  = \frac{2}{3c}$.

\item \emph{All-$W$}: both vertices get $C_{WWW} = 0$.
  Amplitude $= 0$.

\item \emph{Type $(T,T,W)$} (three such assignments):
  each vertex receives one~$W$ half-edge, giving $C_{TTW} = 0$.
  All three vanish.

\item \emph{Type $(T,W,W)$} (three such assignments):
  each vertex receives half-edges $(T,W,W)$, giving
  $C_{TWW} = c$.  The propagator product is
  $\eta^{TT} (\eta^{WW})^2 = (2/c)(9/c^2) = 18/c^3$.
  Each assignment contributes $18c^2/c^3 = 18/c$; with
  $|\mathrm{Aut}| = 12$, the total mixed amplitude is
  \begin{equation}\label{eq:theta-mixed}
  \delta_{\mathrm{theta}} = \frac{3 \cdot 18}{12c}
  = \frac{9}{2c}.
  \end{equation}
\end{itemize}

\subsubsection*{$\Gamma_5$: lollipop}

One self-loop at vertex~$v_0$ (genus~0, valence~3) and one bridge
to vertex~$v_1$ (genus~1, valence~1).  Two edges, four channel
assignments, $|\mathrm{Aut}| = 2$.

\begin{itemize}
\item \emph{$(T,T)$}: self-loop~$T$, bridge~$T$.  Vertex~$v_0$ gets
  half-edges $(T,T,T) \to C_{TTT} = c$.  Vertex~$v_1$ gets
  half-edge $T \to \kappa_T/24 = c/48$.  Propagators:
  $(\eta^{TT})^2 = 4/c^2$.  With $|\mathrm{Aut}|$:
  $\frac{1}{2} \cdot \frac{4}{c^2} \cdot c \cdot \frac{c}{48}
  = \frac{1}{24}$.

\item \emph{$(W,W)$}: self-loop~$W$, bridge~$W$.  Vertex~$v_0$ gets
  $(W,W,W) \to C_{WWW} = 0$.  Vanishes.

\item \emph{$(T,W)$}: self-loop~$T$, bridge~$W$.  Vertex~$v_0$ gets
  $(T,T,W) \to C_{TTW} = 0$.  Vanishes.

\item \emph{$(W,T)$}: self-loop~$W$, bridge~$T$.  Vertex~$v_0$ gets
  $(W,W,T) \to C_{WWT} = c$.  Vertex~$v_1$ gets
  $T \to \kappa_T/24 = c/48$.  Propagators:
  $\eta^{WW} \cdot \eta^{TT} = (3/c)(2/c) = 6/c^2$.
  With $|\mathrm{Aut}|$:
  \begin{equation}\label{eq:lollipop-mixed}
  \delta_{\mathrm{lollipop}}
  = \tfrac{1}{2} \cdot \tfrac{6}{c^2} \cdot c \cdot \tfrac{c}{48}
  = \tfrac{6c^2}{96c^2} = \frac{1}{16}.
  \end{equation}
\end{itemize}
The lollipop mixed amplitude $1/16$ is \emph{independent of~$c$}:
the $c^2$ in the numerator (from $C_{WWT} \cdot \kappa_T/24$)
cancels the $c^2$ in the denominator (from $\eta^{WW} \cdot \eta^{TT}$).


\subsection{The cross-channel correction}
\label{subsec:cross-channel}

Collecting the mixed-channel amplitudes from all boundary graphs:
\begin{equation}\label{eq:cross-channel-total}
\delta F_2
\;=\;
\underbrace{\frac{3}{c}}_{\text{banana}}
\;+\;
\underbrace{\frac{9}{2c}}_{\text{theta}}
\;+\;
\underbrace{\frac{1}{16}}_{\text{lollipop}}
\;+\;
\underbrace{\frac{21}{4c}}_{\text{barbell}}
\;=\;
\frac{51}{4c} + \frac{1}{16}
\;=\;
\frac{c + 204}{16c}.
\end{equation}
The figure-eight and dumbbell graphs contribute zero mixed amplitude
(single edges admit only diagonal assignments).  The correction
decomposes into two structurally distinct pieces:
\begin{itemize}
\item A \emph{$c$-dependent} piece $51/(4c)$, arising from
  genus-0 vertex factors on the banana ($3/c$), theta ($9/(2c)$),
  and barbell ($21/(4c)$) graphs.  This is
  $R$-matrix independent: genus-0 vertices receive no corrections
  from the CohFT $R$-matrix.
\item A \emph{$c$-independent} piece $1/16$, arising from the
  lollipop graph, which mixes a genus-0 vertex factor with a
  genus-1 vertex factor $\kappa_i/24$.
\end{itemize}

\begin{remark}[Comparison with the scalar approximation]
\label{rem:scalar-comparison}
The per-channel (diagonal) graph sum, computed by applying Theorem~D
to each channel independently, gives
\begin{equation}\label{eq:F2-scalar}
F_2^{\mathrm{scal}}(W_3) = \kappa \cdot \lambda_2^{\mathrm{FP}}
= \frac{5c}{6} \cdot \frac{7}{5760} = \frac{7c}{6912}.
\end{equation}
The ratio of the naive cross-channel correction to the scalar
approximation is
\begin{equation}\label{eq:cross-channel-ratio}
\frac{\delta F_2}{F_2^{\mathrm{scal}}}
= \frac{432(c+204)}{7c^2},
\end{equation}
which diverges as $c \to 0$ and decays as $c^{-1}$ at large~$c$.
At $c = 4$: the ratio is approximately~$802$, and the cross-channel
correction \emph{dominates} the scalar term.  At $c = 50$: the
ratio is approximately~$6.3$.
At large~$c$, the lollipop's $c$-independent contribution $1/16$
dominates, and the correction becomes a subleading $O(1/c)$ effect
relative to the $O(c)$ scalar term.
\end{remark}

\begin{remark}[Status of the cross-channel correction]
\label{rem:cross-channel-status}
The computation above uses the naive genus-$1$ vertex factors
$\kappa_i/24$ without $R$-matrix corrections.  The
multi-generator universality problem
(\textbf{op:multi-generator-universality}) is resolved negatively:
the CohFT $R$-matrix does \emph{not} produce compensating terms
that cancel $\delta F_2$.  The cross-channel correction
$\delta F_2 = (c{+}204)/(16c)$ is a genuine additional term,
and the correct decomposition is
$F_g = \kappa \cdot \lambda_g^{\mathrm{FP}} +
\delta F_g^{\mathrm{cross}}$.
\end{remark}


\subsection{The complete amplitude table}
\label{subsec:amplitude-table}

We record the full per-graph data:

\begin{center}
\small
\renewcommand{\arraystretch}{1.3}
\begin{tabular}{lcccc}
\textbf{Graph}
& \textbf{Diagonal ($T$)}
& \textbf{Diagonal ($W$)}
& \textbf{Mixed}
& \textbf{Total} \\
\hline
$\Gamma_1$ figure-eight
& $1/48$ & $1/48$ & $0$ & $1/24$ \\[2pt]
$\Gamma_2$ banana
& $1/c$ & $9/(4c)$ & $3/c$ & $25/(4c)$ \\[2pt]
$\Gamma_3$ dumbbell
& $c/2304$ & $c/3456$ & $0$ & $5c/6912$ \\[2pt]
$\Gamma_4$ theta
& $2/(3c)$ & $0$ & $9/(2c)$ & $31/(6c)$ \\[2pt]
$\Gamma_5$ lollipop
& $1/24$ & $0$ & $1/16$ & $5/48$ \\
\end{tabular}
\end{center}

The table records the boundary-graph amplitudes only; the
smooth graph $\Gamma_0$ contributes an additional integrated
class on $\mathcal{M}_{2,0}$ whose per-channel decomposition
is determined by the constraint
$\sum_\Gamma \mathcal{A}(\Gamma) = \kappa \cdot \lambda_2^{\mathrm{FP}}$
when the full universality identity holds.
The mixed column sums to
$\delta F_2 = (c+204)/(16c)$.

Every entry in this table has been verified against the compute
module \texttt{w3\_genus2.py}, using exact rational arithmetic
(\texttt{fractions.Fraction}), at the seven test central charges
$c \in \{1, 2, 4, 13, 26, 50, 100\}$.


\subsection{Planted-forest corrections}
\label{subsec:planted-forest}

The planted-forest correction on each primary line is
$\delta_{\mathrm{pf}}^{(2,0)} = S_3(10 S_3 - \kappa)/48$.

\emph{$T$-line}: $S_3^T = 2$ (the Virasoro cubic shadow),
$\kappa_T = c/2$.
\begin{equation}\label{eq:pf-T}
\delta_{\mathrm{pf}}^T
= \frac{2(20 - c/2)}{48}
= \frac{40 - c}{48}.
\end{equation}
This vanishes at $c = 40$, is positive for $c < 40$, and negative
for $c > 40$.

\emph{$W$-line}: $S_3^W = 0$ by the $\mathbb{Z}_2$ parity
$W \mapsto -W$, which forces all odd-arity shadows on the $W$-line
to vanish.  Hence $\delta_{\mathrm{pf}}^W = 0$.

The total planted-forest correction is $\delta_{\mathrm{pf}} =
(40-c)/48$, living entirely on the $T$-line.


\subsection{Koszul duality constraints}
\label{subsec:w3-koszul-duality}

Under the Koszul involution, $W_3$ at central charge~$c$ is
dual to $W_3$ at central charge $c' = 100 - c$
(the Feigin--Frenkel involution for $\mathfrak{sl}_3$).  The modular
characteristics satisfy
\begin{equation}
\kappa(c) + \kappa(100-c) = \frac{5c}{6} + \frac{5(100-c)}{6}
= \frac{250}{3},
\end{equation}
and the scalar free energies satisfy $F_2^{\mathrm{scal}}(c) +
F_2^{\mathrm{scal}}(100-c) = (250/3) \cdot \lambda_2^{\mathrm{FP}}$.
The cross-channel correction at the dual pair is
\[
\delta F_2(c) + \delta F_2(100-c)
= \frac{c+204}{16c} + \frac{304-c}{16(100-c)}
= -\,\frac{c^2 - 100c - 10200}{8c(100-c)}.
\]
This sum is not constant in~$c$: the cross-channel correction does
not respect Koszul additivity.  If universality holds ($\delta F_2$
is cancelled by $R$-matrix corrections), this is automatic.  If
the naive correction persists, it produces a Koszul-asymmetric
contribution to $F_2$---an observable departure from single-channel
behaviour.


%======================================================================
\section{The gravitational coproduct is primitive}
\label{sec:gravitational-coproduct}
%======================================================================

In gauge theory, the gluon is a fundamental field: it can split.
A single gluon entering a scattering region can exit as two gluons,
and the coproduct of the dg-shifted Yangian encodes this splitting
at the algebraic level.  In gravity, the graviton is not fundamental
but composite.  The stress tensor is the Sugawara quadratic
\begin{equation}\label{eq:sugawara}
T = \frac{1}{k+2}\Bigl(\tfrac{1}{2} {:}J^0 J^0{:}
+ {:}J^+ J^-{:}\Bigr),
\end{equation}
built from two copies of the affine connection by normal ordering.
The compositeness of the graviton has a precise algebraic
consequence: the transferred coproduct is trivial.


\subsection{The DS-HPL transfer theorem}
\label{subsec:ds-hpl}

The Drinfeld--Sokolov reduction $\hat{\mathfrak{sl}}_2 \to
\mathrm{Vir}$ is a BRST reduction with a strong deformation retract
$(i, p, h)$ from the BRST complex to the Virasoro cohomology.
The homotopy $h$ has ghost-number degree~$-1$: it maps
$V \mapsto b$ (ghost number $0 \to -1$) and
$c \mapsto U/2$ (ghost number $1 \to 0$).

\begin{theorem}[DS-HPL transfer]\label{thm:ds-hpl-standalone}
Let $\fg$ be a simple Lie algebra, and let
$Y^{\mathrm{dg}}(\widehat{\fg})$ denote the dg-shifted
Yangian at level~$k$, with $A_\infty$ products
$\{m_n^{\mathrm{aff}}\}$, coproduct
$\Delta_z^{\mathrm{aff}}$, and $r$-matrix
$r^{\mathrm{aff}}(z) = \Omega/z$.  The HPL transfer through the
principal BRST SDR $(i, p, h)$ from the DS complex to
$H^*_{\mathrm{BRST}} = \cW(\fg)$ produces:
\begin{enumerate}[label=\textup{(\roman*)}]
\item Transferred $A_\infty$ products
  $\{m_n^{\cW}\}_{n \geq 2}$ on the $\cW$-algebra sector,
  with $m_n^{\cW} \neq 0$ for all $n \geq 2$
  \textup{(}the entire infinite $A_\infty$ tower of class~$\mathbf{M}$%
  \textup{)}.
\item A transferred coproduct
  $\Delta_z^{\cW} \colon H_{\cW} \to
  H_{\cW} \otimes H_{\cW}$.
\item A transferred $r$-matrix
  $r^{\cW}(z)$ with leading pole of order
  $2h_{\max} - 1$, where $h_{\max}$ is the highest conformal
  weight among the generators of~$\cW(\fg)$.
  For $\fg = \mathfrak{sl}_2$:
  $r^{\mathrm{Vir}}(z) = (c/2)/z^3 + 2T/z$.
\end{enumerate}
\end{theorem}

The $r$-matrix deserves comment.  The affine $r$-matrix
$r^{\mathrm{aff}}(z) = \Omega/z$ is a simple pole---the
collision residue of the affine bar complex, whose OPE has at most
a double pole.  The Virasoro OPE has a quartic pole
($z^{-4}$ in $T(z)T(w)$); the bar kernel $d\!\log(z-w)$ absorbs
one power, producing the cubic pole $z^{-3}$ of
$r^{\mathrm{Vir}}(z)$.  This is the first instance of the general
principle: the $r$-matrix pole order equals the OPE pole order
minus one.

The key result is the following.

\begin{theorem}[Gravitational coproduct primitivity]
\label{thm:grav-primitivity-standalone}
For any simple Lie algebra~$\fg$, the HPL-transferred coproduct
of the principal $\cW$-algebra $\cW(\fg)$ is strictly primitive
at all arities:
\begin{equation}\label{eq:grav-prim}
\Delta_{z,k}^{\cW(\fg)} = 0
\qquad \text{for all } k \geq 2.
\end{equation}
In particular, for $\fg = \mathfrak{sl}_2$:
$\Delta_z^{\mathrm{Vir}}(T) = \tau_z(T) \otimes 1
+ 1 \otimes T$.  All nontrivial structure of the $\cW$-algebra
dg-shifted Yangian resides in the $r(z)$-twisted target product,
never in the coproduct.
\end{theorem}

\begin{proof}
The proof is a ghost-number obstruction.  The argument is
uniform in~$\fg$: only the ghost-number grading of the
principal BRST complex is used, and this grading is universal
for all principal DS reductions.  We write the proof for
general~$\fg$, illustrating with
$\fg = \mathfrak{sl}_2$ where explicit.

\emph{Step~1 (linear primitivity).}
The linear transferred coproduct
$(p \otimes p)(\Delta_z(\phi))$ is already primitive for
every generator~$\phi$ of $\cW(\fg)$.  The affine coproduct
$\Delta_z(\phi)$ involves the Casimir element
$\Omega = \sum_a J^a \otimes J_a$, and the projection~$p$
onto the $\cW$-algebra cohomology $H^*_{\mathrm{BRST}}$
annihilates the affine currents~$J^a$ (which have conformal
weight~$1$ and ghost number~$0$, but lie outside the
$\cW$-algebra sector).  Since every summand of~$\Omega$
contains at least one $J^a$ factor, $(p \otimes p)(\Omega) = 0$.

For $\fg = \mathfrak{sl}_2$: $p$ annihilates $J^0$, $J^\pm$;
the Casimir $\Omega(w{+}z,w) = J^0 \otimes J^0
+ J^+ \otimes J^- + J^- \otimes J^+$ is killed by
$p \otimes p$.

\emph{Step~2 (ghost-number obstruction at higher arity).}
At arity $k \geq 2$, every HPL tree has at least one internal
edge labelled by the homotopy~$h$.  The ghost-number accounting:
\begin{center}
\small
\begin{tabular}{lc}
Operation & Ghost degree \\
\hline
$h$ (homotopy) & $-1$ \\
$\delta$ (BRST perturbation) & $+1$ \\
$\mu$ (vertex algebra product) & $0$ \\
$\Delta_z$ (coproduct) & $0$ \\
\end{tabular}
\end{center}
The product $h\delta$ has ghost degree $-1 + 1 = 0$, but the
tree topology at arity $k \geq 2$ always contains one
\emph{uncompensated}~$h$: the outermost homotopy that is not
preceded by a $\delta$ insertion.  The overall output before
projection has ghost number $\leq -1$.  This counting is
independent of~$\fg$: it depends only on the tree
structure of the HPL formula and the universal ghost-number
assignments.

\emph{Step~3 (projection kills everything).}
The $\cW$-algebra cohomology sits in ghost number~$0$.
The projection $p$ annihilates all ghost-number-$(\neq 0)$ elements.
Since the output of each arity-$k$ tree has ghost number $\leq -1$,
the projection $(p \otimes p)$ kills every summand.

\emph{Step~4 (perturbative stability).}
The BRST perturbation $\delta$ has ghost degree~$+1$, so the
product $h\delta$ preserves ghost number.  For every generator
$\phi \in \cW(\fg)$: $\delta(\phi) = 0$ ($\cW$-generators are
BRST-closed by construction of the DS reduction), so
$(h\delta)(i(\phi)) = 0$, and by induction
$(h\delta)^n(i(\phi)) = 0$ for all $n \geq 1$.  The corrected
inclusion $i_\infty(\phi) = i(\phi)$ is unperturbed, and the
same primitive coproduct results at every perturbative order.
\end{proof}


\subsection{Physical meaning}
\label{subsec:physical-meaning}

The coproduct $\Delta_z$ of a dg-shifted Yangian encodes the
\emph{splitting} of excitations.  Its strictness or nontriviality
is the algebraic fingerprint of the distinction between fundamental
and composite fields.

\begin{center}
\small
\renewcommand{\arraystretch}{1.3}
\begin{tabular}{llccl}
\textbf{Theory}
& \textbf{Boundary algebra}
& $\boldsymbol{\Delta_z}$
& \textbf{Shadow depth}
& \textbf{Reason} \\
\hline
Chern--Simons &
  $\hat{\mathfrak{g}}_k$ (affine KM) &
  $\neq \varepsilon$ &
  $3$ (class $\mathbf{L}$) &
  fundamental gluons split \\[4pt]
3d gravity &
  $\mathrm{Vir}_c = W_k(\mathfrak{sl}_2)$ &
  $= \varepsilon$ &
  $\infty$ (class $\mathbf{M}$) &
  ghost obstruction (DS) \\[4pt]
Twisted holography &
  $H_k$ (Heisenberg) &
  $= \varepsilon$ &
  $2$ (class $\mathbf{G}$) &
  free field, $r = k/z$ \\[4pt]
M2 brane &
  DDCA &
  $\neq \varepsilon$ &
  --- &
  no DS, fundamental modes \\[4pt]
M5 brane &
  $W_{1+\infty}(K)$ &
  $= \varepsilon$ &
  $\infty$ (class $\mathbf{M}$) &
  ghost obstruction (DS) \\
\end{tabular}
\end{center}

The pattern is clean: \emph{theories arising from principal
Drinfeld--Sokolov reduction have primitive coproducts}.  The
ghost sector of the BRST complex provides the grading obstruction
that kills all higher-arity coproduct components.
Theorem~\ref{thm:grav-primitivity-standalone} proves this for
all principal DS reductions $\widehat{\fg} \to \cW(\fg)$: the
ghost-number grading of the BRST complex is universal.
Theories with fundamental excitations (Chern--Simons, deformed
double current algebras) have no ghost sector available, and
their coproducts are nontrivial.

The two exceptions to the DS pattern are illuminating.  The
Heisenberg algebra has a primitive coproduct for a different reason:
it is a free field, its $r$-matrix $r^H(z) = k/z$ is scalar, and
the bar complex is formal at arity~$2$ (shadow depth $r_{\max} = 2$,
class~$\mathbf{G}$).  No DS reduction and no ghost obstruction;
primitivity follows from the absence of nonlinear OPE structure.
The M2 brane boundary algebra (a deformed double current algebra)
has a nontrivial coproduct because it is not obtained by DS
reduction from any affine algebra---its excitations are fundamental
modes of the M2 worldvolume theory.


\subsection{The \texorpdfstring{$r$}{r}-matrix carries all
gravitational structure}
\label{subsec:r-matrix-gravity}

With the coproduct primitive, the full gravitational content of
the Virasoro dg-shifted Yangian is encoded in the $r$-matrix
\begin{equation}\label{eq:grav-r-matrix}
r^{\mathrm{Vir}}(z) = \frac{c/2}{z^3} + \frac{2T}{z}.
\end{equation}
This two-term expression carries a remarkable amount of structure:
\begin{itemize}
\item The \emph{leading coefficient} $c/2 = \kappa(\mathrm{Vir}_c)$
  is the modular characteristic.  It determines the genus tower
  $F_g = (c/2) \int_{\overline{\mathcal{M}}_g} \lambda_g$ at all
  genera, the BTZ black hole entropy
  $S_{\mathrm{BTZ}} = 2\pi\sqrt{ch/6}$, and the Hawking--Page
  transition at $h = c/24$.

\item The \emph{subleading term} $2T/z$ encodes the state dependence
  of gravitational scattering.  The braiding of Wilson lines
  (gravitational line operators) in the holomorphic-topological
  bulk is controlled by this term.

\item The \emph{pole structure}---cubic leading pole, simple
  subleading pole, no even poles---is the signature of a
  bosonic algebra extracted through the $d\!\log$ kernel: the
  OPE poles $z^{-4}, z^{-2}, z^{-1}$ become $r$-matrix poles
  $z^{-3}, z^{-1}$ (each shifted by $-1$), with the $z^{-2}$
  OPE pole producing $z^{-1}$ (which merges with the subleading
  term from $z^{-1}$).
\end{itemize}

The contrast with Chern--Simons is maximal.  In Chern--Simons
gauge theory, the affine $r$-matrix $r^{\mathrm{aff}}(z) = \Omega/z$
is a simple pole and the coproduct is nontrivial: gluons split
and the splitting is nontrivial.  In gravity, the $r$-matrix has
a cubic pole and the coproduct is trivial: gravitons braid but
do not split.  The Drinfeld--Sokolov reduction bridges the two:
it takes the affine Casimir $\Omega/z$ to the gravitational
$(c/2)/z^3 + 2T/z$ by homological perturbation through the
Sugawara nonlinearity, and in the same stroke, the ghost-number
obstruction kills the nontrivial coproduct.

\begin{center}
\small
\renewcommand{\arraystretch}{1.15}
\begin{tabular}{lll}
\textbf{Datum} & \textbf{Chern--Simons} & \textbf{Gravity} \\
\hline
boundary algebra
  & $V_k(\mathfrak{g})$ (affine KM)
  & $\mathrm{Vir}_c$ (Virasoro) \\[2pt]
highest OPE pole
  & $z^{-2}$ (quadratic)
  & $z^{-4}$ (quartic) \\[2pt]
$A_\infty$ tower
  & strict ($m_{k \geq 3} = 0$)
  & infinite ($m_k \neq 0$ all $k$) \\[2pt]
shadow depth
  & $3$ (class $\mathbf{L}$)
  & $\infty$ (class $\mathbf{M}$) \\[2pt]
$r^{\mathrm{coll}}(z)$
  & $\Omega/z$ (simple pole)
  & $(c/2)/z^3 + 2T/z$ (cubic pole) \\[2pt]
$\Delta_z$
  & nontrivial (gluons split)
  & primitive (gravitons braid) \\[2pt]
$\kappa$
  & $\dim\mathfrak{g}\,(k+h^\vee)/(2h^\vee)$
  & $c/2$ \\[2pt]
$\kappa + \kappa^!$
  & $0$
  & $13$
\end{tabular}
\end{center}

The entire table descends from a single datum: the number of
$\lambda$-powers in the $\lambda$-bracket.  Affine Kac--Moody has
$\lambda^1$ at most (quadratic pole), yielding strict associativity,
finite shadow depth, a simple-pole $r$-matrix, and a nontrivial
coproduct.  The Virasoro algebra has $\lambda^3$ (quartic pole),
yielding the full infinite $A_\infty$ tower, infinite shadow depth,
a cubic-pole $r$-matrix, and a primitive coproduct.  The
Drinfeld--Sokolov reduction is the functor that connects the two
columns: it increases the $\lambda$-bracket order by two, trades
the nontrivial coproduct for a ghost obstruction, and replaces the
complementarity $\kappa + \kappa^! = 0$ with $\kappa + \kappa^! = 13$.

\smallskip

In one sentence:
\emph{gravity does not produce particles, only braids; the graviton
is composite, and the ghost sector of its composite origin is the
algebraic obstruction that kills the splitting}.
